% ═══════════════════════════════════════════════════════════════
% AB-08a: Viajar -- Reisen und Verkehrsmittel
% Abgeleitet aus LE-08a (Score 9.3/10)
% ═══════════════════════════════════════════════════════════════

\documentclass[11pt, a4paper]{article}

% ═══════════════════════════════════════════════════════════════
% SW-MODUS TOGGLE
% ═══════════════════════════════════════════════════════════════
% Setze \swmodustrue für Schwarz-Weiß-Druck
\newif\ifswmodus
\swmodusfalse    % Standard: Farbe

\usepackage[utf8]{inputenc}
\usepackage[T1]{fontenc}
\usepackage[ngerman]{babel}
\usepackage{helvet}
\renewcommand{\familydefault}{\sfdefault}

\usepackage[margin=2cm]{geometry}
\usepackage{enumitem}
\usepackage{tabularx}
\usepackage{booktabs}

\usepackage{xcolor}
\usepackage{tcolorbox}
\tcbuselibrary{skins}

\usepackage{fancyhdr}
\usepackage{lastpage}
\usepackage{needspace}          % Für \needspace (Seitenumbruch-Kontrolle)

% Farben - Basis
\definecolor{spanisch-color}{HTML}{c41e3a}
\definecolor{grau-color}{HTML}{888888}
\definecolor{hellgrau-color}{HTML}{f5f5f5}
\definecolor{orange-color}{HTML}{b35c00}

% SW-Farben (druckerfreundlich: keine Füllflächen, nur Linien)
\definecolor{sw-dark}{HTML}{000000}
\definecolor{sw-gray}{HTML}{666666}
\definecolor{sw-white}{HTML}{ffffff}

% Bedingte Farbzuweisung
\ifswmodus
    \colorlet{spanisch}{sw-dark}       % Rahmen/Text: schwarz
    \colorlet{grau}{sw-gray}           % Sekundärtext: dunkelgrau
    \colorlet{hellgrau}{sw-white}      % Hintergrund: weiß (keine Füllung)
    \colorlet{orange}{sw-dark}         % Hinweis-Rahmen: schwarz
\else
    \colorlet{spanisch}{spanisch-color}
    \colorlet{grau}{grau-color}
    \colorlet{hellgrau}{hellgrau-color}
    \colorlet{orange}{orange-color}
\fi

\newcommand{\fachfarbe}{spanisch}

% Variablen
\newcommand{\thema}{Viajar -- Reisen und Verkehrsmittel}
\newcommand{\sequenz}{08}
\newcommand{\block}{a}
\newcommand{\fach}{Spanisch}
\newcommand{\klasse}{7}
\newcommand{\maxpunkte}{24}

% Kopf-/Fußzeile
\pagestyle{fancy}
\fancyhf{}
\renewcommand{\headrulewidth}{0.4pt}
\renewcommand{\footrulewidth}{0.4pt}
\lhead{\fach{} -- Klasse \klasse}
\chead{AB-\sequenz\block}
\rhead{}
\lfoot{}
\rfoot{Seite \thepage\ von \pageref{LastPage}}

% Befehle
\newcommand{\punktebox}[1]{\hfill\fbox{\small \_/#1}}
\newcommand{\luecke}[1]{\rule{#1}{0.4pt}}
\newcommand{\es}[1]{\textcolor{\fachfarbe}{\textbf{#1}}}

% Merkkasten (SW-kompatibel: kein gefüllter Titelbalken)
\newtcolorbox{merkkasten}[1][]{
    colback=hellgrau,
    colframe=\fachfarbe,
    colbacktitle=white,
    coltitle=\fachfarbe,
    boxrule=1pt,
    arc=0pt,
    left=8pt, right=8pt, top=8pt, bottom=8pt,
    fonttitle=\bfseries,
    attach boxed title to top left={yshift=-2mm, xshift=5mm},
    boxed title style={boxrule=0pt, colback=white},
    title=#1
}

% Hinweis (SW-kompatibel: kein gefüllter Titelbalken)
\newtcolorbox{hinweis}{
    colback=white,
    colframe=orange,
    colbacktitle=white,
    coltitle=orange,
    boxrule=1pt,
    arc=0pt,
    left=5pt, right=5pt, top=5pt, bottom=5pt,
    fonttitle=\bfseries,
    attach boxed title to top left={yshift=-2mm, xshift=5mm},
    boxed title style={boxrule=0pt, colback=white},
    title=Hinweis
}

% Aufgabe (kein Seitenumbruch innerhalb)
\newenvironment{aufgabe}[1]{%
    \needspace{5\baselineskip}%     % Mindestens 5 Zeilen Platz, sonst neue Seite
    \nopagebreak[4]%
    \vspace{10pt}
    \noindent\textbf{\textcolor{\fachfarbe}{Aufgabe #1}}
    \nopagebreak[4]%
    \vspace{5pt}
    \nopagebreak[4]%
    \begin{list}{}{\leftmargin=0pt}
    \item[]
}{%
    \end{list}
}

% ═══════════════════════════════════════════════════════════════
\begin{document}

% Kopfbereich
\begin{center}
    {\Large\bfseries\textcolor{\fachfarbe}{Arbeitsblatt: \thema}}
\end{center}

\vspace{5pt}

\noindent
\begin{tabularx}{\textwidth}{|X|X|X|}
\hline
\textbf{Name:} \luecke{4cm} &
\textbf{Klasse:} \klasse &
\textbf{Datum:} \luecke{2cm} \\
\hline
\end{tabularx}

\vspace{15pt}

% ═══════════════════════════════════════════════════════════════
% MERKKASTEN 1: Verkehrsmittel ($\rightarrow$ FZ1)
% ═══════════════════════════════════════════════════════════════

\begin{merkkasten}[Merkkasten: Los medios de transporte (Verkehrsmittel)]

\begin{center}
\begin{tabular}{|l|l|l|}
\hline
\textbf{Spanisch} & \textbf{Deutsch} & \textbf{Beispiel} \\
\hline
\luecke{3cm} & das Flugzeug & Vamos \luecke{1.5cm} avión. \\
\hline
\luecke{3cm} & der Zug & El tren sale a las ocho. \\
\hline
\luecke{3cm} & das Auto & Viajamos en coche. \\
\hline
\luecke{3cm} & der Bus & Tomo el autobús. \\
\hline
\luecke{3cm} & das Schiff & El barco va a las islas. \\
\hline
\end{tabular}
\end{center}

\vspace{0.5em}
\textit{Übertrage die Vokabeln aus dem Lernpfad oder von der Tafel!}
\end{merkkasten}

\vspace{10pt}

% ═══════════════════════════════════════════════════════════════
% MERKKASTEN 2: ir a + Infinitiv ($\rightarrow$ FZ3)
% ═══════════════════════════════════════════════════════════════

\begin{merkkasten}[Merkkasten: ir a + Infinitiv (nahe Zukunft)]

\begin{center}
\begin{tabular}{|l|l|l|}
\hline
\textbf{Person} & \textbf{ir} & \textbf{+ a + Infinitiv} \\
\hline
yo & \luecke{2cm} & a nadar \\
\hline
tú & \luecke{2cm} & a viajar \\
\hline
él/ella & \luecke{2cm} & a hacer senderismo \\
\hline
nosotros & \luecke{2cm} & a tomar el sol \\
\hline
\end{tabular}
\end{center}

\vspace{0.5em}
\textbf{Bedeutung:} \textit{Voy a nadar} = Ich werde schwimmen / Ich gehe schwimmen.
\end{merkkasten}

\vspace{15pt}

% ═══════════════════════════════════════════════════════════════
% AUFGABE 1: Verkehrsmittel benennen ($\rightarrow$ FZ1, AFB I)
% ═══════════════════════════════════════════════════════════════

\begin{aufgabe}{1}
Schreibe die spanischen Wörter für die Verkehrsmittel. \punktebox{5}
\end{aufgabe}

\begin{hinweis}
\textbf{Wortbank:} el avión | el tren | el coche | el autobús | el barco
\end{hinweis}

\vspace{0.5em}

\begin{enumerate}[label=\alph*)]
    \item Flugzeug: \luecke{4cm}
    \item Zug: \luecke{4cm}
    \item Auto: \luecke{4cm}
    \item Bus: \luecke{4cm}
    \item Schiff: \luecke{4cm}
\end{enumerate}

\vspace{15pt}

% ═══════════════════════════════════════════════════════════════
% AUFGABE 2: Zuordnung Orte + Aktivitäten ($\rightarrow$ FZ2, AFB I)
% ═══════════════════════════════════════════════════════════════

\begin{aufgabe}{2}
Verbinde die Orte mit den passenden Aktivitäten. \punktebox{4}
\end{aufgabe}

\begin{center}
\begin{tabular}{rcl}
la playa & \luecke{3cm} & hacer senderismo \\[0.8em]
la montaña & \luecke{3cm} & nadar, tomar el sol \\[0.8em]
el hotel & \luecke{3cm} & hacer una barbacoa \\[0.8em]
el camping & \luecke{3cm} & descansar \\
\end{tabular}
\end{center}

\vspace{15pt}

% ═══════════════════════════════════════════════════════════════
% AUFGABE 3: Sätze mit ir a + Infinitiv ($\rightarrow$ FZ3, AFB II)
% ═══════════════════════════════════════════════════════════════

\begin{aufgabe}{3}
Ergänze die Sätze mit der richtigen Form von \textit{ir a + Infinitiv}. \punktebox{6}
\end{aufgabe}

\begin{hinweis}
\textbf{Verben:} nadar | viajar | hacer | tomar | ir
\end{hinweis}

\vspace{0.5em}

\begin{enumerate}[label=\alph*)]
    \item Yo \luecke{4cm} en la playa. \textit{(schwimmen)}
    \item Nosotros \luecke{4cm} en avión. \textit{(reisen)}
    \item Tú \luecke{4cm} senderismo. \textit{(machen)}
    \item Ellos \luecke{4cm} el sol. \textit{(nehmen/sich sonnen)}
    \item Ella \luecke{4cm} a la montaña. \textit{(gehen)}
    \item Vosotros \luecke{4cm} en el hotel. \textit{(bleiben/quedarse)}
\end{enumerate}

\vspace{15pt}

% ═══════════════════════════════════════════════════════════════
% AUFGABE 4: Dialog über Urlaubspläne ($\rightarrow$ FZ4, AFB II)
% ═══════════════════════════════════════════════════════════════

\begin{aufgabe}{4}
Vervollständige den Dialog. Nutze die Redemittel. \punktebox{6}
\end{aufgabe}

\begin{hinweis}
\textbf{Redemittel:} ¿Adónde vas a ir? | Voy a ir a... | ¿Cómo vas a viajar? | Voy a viajar en... | ¿Qué vas a hacer? | Voy a...
\end{hinweis}

\vspace{0.5em}

\noindent
\textbf{A:} \es{¿Adónde vas a ir de vacaciones?}

\textbf{B:} \luecke{8cm}

\textbf{A:} \es{¿Cómo vas a viajar?}

\textbf{B:} \luecke{8cm}

\textbf{A:} \luecke{6cm}

\textbf{B:} \luecke{8cm}

\vspace{15pt}

% ═══════════════════════════════════════════════════════════════
% AUFGABE 5: Reiseplan schreiben ($\rightarrow$ FZ5, AFB III)
% ═══════════════════════════════════════════════════════════════

\begin{aufgabe}{5}
Schreibe einen kurzen Reiseplan (mind. 5 Sätze). \punktebox{3}
\end{aufgabe}

\begin{hinweis}
Schreibe über: Wohin? Mit welchem Verkehrsmittel? Was machst du dort?\\
\textbf{Beispiel:} \textit{Este verano voy a ir a España. Voy a viajar en avión...}
\end{hinweis}

\noindent\rule{\textwidth}{0.4pt}\vspace{8pt}
\noindent\rule{\textwidth}{0.4pt}\vspace{8pt}
\noindent\rule{\textwidth}{0.4pt}\vspace{8pt}
\noindent\rule{\textwidth}{0.4pt}\vspace{8pt}
\noindent\rule{\textwidth}{0.4pt}\vspace{8pt}

\vspace{20pt}
\hrule
\vspace{10pt}

\begin{flushright}
\textbf{Punkte:} \luecke{1cm} / \maxpunkte
\end{flushright}

\end{document}
